\section{Увод}
\label{sec:introduction}

\subsection{Постановка на проблема}

Достъпът до данни остава едно от местата, в които разминаването между езика на приложението и модела на съхранение поражда най-много структурни грешки. В традиционните слоеве за достъп до бази данни заявките се описват като низове, които се проверяват едва при изпълнение (runtime). Това означава, че неправилни имена на колони, несъвместими типове, невалидни конструкции за свързване (\code{JOIN}) или неправилно подадени параметри се откриват късно --- често върху реална база данни и след допълнителни ръчни тестове. Същевременно смяната на системата за съхранение (backend) от релационна към документна, графова или ключ-стойност (key-value) обикновено изисква пълно пренаписване на слоя за достъп до данни.

Обектно-релационното картографиране (Object-Relational Mapping, ORM) е класическият подход за свързване на обектния модел на приложението с персистентното съхранение~\cite{fowler2002}. Повечето реализации обаче остават силно зависими от механизми, действащи по време на изпълнение: текстови заявки на структурирания език за заявки (Structured Query Language, SQL), кодогенерация, анотации, макроси или динамично описвани схеми. В контекста на C++ това положение е особено противоречиво. Езикът предлага мощна шаблонна система, изчислимост по време на компилация (compile-time) чрез \code{constexpr} и концепции (concepts), но в библиотеките за достъп до данни тези механизми често не се използват в пълна степен. Настоящата работа изследва доколко е възможно компилаторът да поеме ролята на първи валидатор на структурата на заявката и на съвместимостта с конкретната система за съхранение.

Съществена особеност на разработваната библиотека е, че въпреки името „обектно-релационно картографиране“, нейната архитектура не е ограничена само до релационния модел. Напротив, тя използва един и същ C++ модел на данните и едно и също типизирано междинно представяне на заявката (query intermediate representation, query IR) като логически слой над релационни и нерелационни системи. Това поставя втория изследователски въпрос: кои части на ORM абстракцията са универсални и кои неизбежно трябва да се ограничават чрез договор, базиран на способности (capability-based contract).

Третият проблемен пласт е свързан с начина на изпълнение. Дори когато логическият модел и междинното представяне са типово коректни, практическата стойност на една съвременна библиотека за достъп до данни зависи и от това дали тя интегрира асинхронно изпълнение, без да разруши типовите гаранции и без принудително преминаване към стил, основан на обратни извиквания (callbacks). Така дипломната работа се организира около две допълващи се оси: моделиране по време на компилация и асинхронна архитектура на изпълнението, изградена върху корутини (coroutines).

\subsection{Цел и задачи на разработката}

Основната цел на дипломната работа е да представи проектирането и реализацията на библиотека за обектно-релационно картографиране по време на компилация за C++, която:

\begin{enumerate}[label=\arabic*.]
    \item моделира данните и заявките чрез статично типизирани C++ конструкции;
    \item позволява изграждане на \code{SELECT}, \code{INSERT}, \code{UPDATE} и \code{DELETE} заявки като \code{constexpr} обекти;
    \item отделя логическата структура на заявката от специфичната за системата за съхранение материализация;
    \item поддържа релационни и нерелационни конектори чрез \code{connector\_trait<DB>};
    \item използва механизъм на способностите за ограничаване на неподдържани операции още при компилация;
    \item демонстрира как C++ моделът определя логическата схема и физическата структура на данните при различни системи за съхранение;
    \item реализира асинхронен слой, основан на корутини, който надгражда синхронния модел по формален и типово безопасен начин;
    \item формализира договор за добавяне на нови конектори, включително такива с неблокираща поддръжка от страна на драйвера;
    \item валидира архитектурата чрез систематични единични (unit) и интеграционни тестове и емпирична оценка.
\end{enumerate}

За постигането на тази цел работата решава следните задачи: анализ на съществуващите решения; дефиниране на статичен модел на обектите; изграждане на междинно представяне на заявката; проектиране на конекторния слой и маркерите за способности; изграждане на асинхронен слой, основан на корутини; реализация на конектори за релационни и нерелационни бази данни; изследване на миграции, подготвени заявки, нишкова безопасност и асинхронно управление на връзките; формулиране на критерии за разширяемост; изграждане на методология за емпирична оценка на синхронния и асинхронния режим.

\subsection{Обхват и ограничения}

Работата разглежда библиотеката като слой за моделиране, изграждане и изпълнение на заявки. В обхвата ѝ попадат: моделът на обектите, междинното представяне на заявката, механизмът на заместителите (placeholders), подготвените заявки, миграционният модел, архитектурата на конекторния слой, ограниченията чрез маркери за способности, асинхронният слой на корутините и конкретните реализации на конектори. Емпиричната оценка покрива релационни (SQLite, PostgreSQL, MySQL), документни (MongoDB) и ключ-стойност (Redis) системи; ширококолонният (Cassandra) и графовият (Neo4j) сценарий се разглеждат само в обобщителен план поради ограниченията на обема.

Извън непосредствения обхват остават оптимизации на производствено равнище като разпределено управление на връзките, автоматизирани миграционни стратегии за всички системи за съхранение, оценка върху пълна матрица от производствени среди и пълно използване на бъдещата функционалност за отражение (reflection) на C++26. Тези направления са разгледани като перспектива в Глава~\ref{sec:conclusion}.

\subsection{Структура на изложението}

Дипломната работа е организирана по следния начин. Глава~\ref{sec:existing_solutions} разглежда съществуващите C++ решения за достъп до бази данни и подходите за обектно-релационно картографиране и позиционира настоящата разработка спрямо тях. Глава~\ref{sec:design} представя проектните решения: теоретичната основа на обектно-релационното картографиране, моделите за съхранение, слоестата организация на библиотеката, обектния и заявковия слой, формалния договор на конектора, механизма на способностите и теорията на корутините и асинхронния дизайн. Глава~\ref{sec:implementation} описва програмната реализация: асинхронния изпълнителен слой, конкретните конектори за основните системи за съхранение и подробната верификация чрез единични, интеграционни и емпирични тестове. Глава~\ref{sec:conclusion} обобщава приносите, ограниченията и направленията за бъдещо развитие.
